\section{The chain: a real corner, that lasts, that predicts}
\label{sec:synthesis}

This chapter is the keystone of the third pillar: it shows that the three payoffs of operating in the latent
space are not three lucky findings but one mechanism. The earlier chapters each answered a question in
isolation --- is the space a continuum or hidden biotypes (\S\ref{sec:strata}), does it hold still over time
(\S\ref{sec:temporal}), does it forecast the future (\S\ref{sec:prognosis}). Read separately they are three
results. Read together they are a single argument, and the argument is worth slowing down for, because its
logic is the spine of the whole program: \emph{the part of the map that does not move is the part that tells
you where the patient is going}. This short chapter is a tutorial on that one chain. It introduces no new
model and no new number; it explains how three findings lock into each other, so that the conclusion --- a
biological corner of a continuum, durable in time, predictive of functioning --- reads as one claim rather
than three coincidences, and so that the clinical value of the latent-space view (stratify on the durable
biology, forecast functioning, monitor the moving symptoms) reads as \emph{earned}.

\begin{figure}[h]\centering
\includegraphics[width=0.96\linewidth]{figures/synthesis.png}
\caption{The chain in one picture. \emph{Left (M2):} the transdiagnostic space is a graded continuum, not
a set of biotypes; its most distinct extreme is the \emph{biological} corner (archetype $A_0$: high
metabolic, inflammatory, substance load). \emph{Middle (M3):} scored forward onto follow-up without
re-estimation, that biological corner is the \emph{durable} part of the map --- metabolic test--retest
$\text{ICC}=0.91$, a trait --- while symptoms slide. \emph{Right (M4):} that same durable biological corner
forecasts two-year \emph{functioning}; the patients sitting in it have the worst functional-remission
prospect ($27\%$, against $60\%$ at the low-burden pole). One arrow, three chapters: \emph{a real corner,
that lasts, that predicts}.}
\label{fig:synthesis}
\end{figure}

\subsection{Link one --- M2: a real continuum, with biology as its most distinct corner}
The first link establishes the \emph{geometry} the rest of the chain stands on. The coordinate-system chapter
(\S\ref{sec:strata}) asked the field's oldest question of these data --- are there biotypes? --- and
answered \emph{no}, by trying hard to find them and failing on every test. A single-Gaussian falsification
null is the cleanest way to see this: the best clustering of the real coordinates scores a silhouette of
$0.140$, and a smooth blob simulated to have \emph{no} clusters at all scores $0.141\pm0.002$ (a separation
of $z=-0.36$). The data are statistically indistinguishable from a structureless cloud as far as cluster
boundaries are concerned. Every test that probes for \emph{separation} concurs: the principal axis is unimodal
(dip $p\approx1$), density clustering finds nothing (HDBSCAN returns zero clusters), the topology is a single
connected component, and once each patient's measurement uncertainty is propagated the optimal mixture
collapses to one component in all twenty posterior draws. The two tests that \emph{do} fire --- Hopkins
clusterability and a GMM-BIC interior optimum --- are precisely the clustering-\emph{tendency} signals the null
reproduces, so they register the cloud's non-Gaussian \emph{shape}, not a set of gaps.

\begin{definitionbox}[Continuum, not biotypes]
``Continuum'' is not a hedge; it is a finding with teeth. It says there are no natural cut-points --- no
place where one patient group ends and another begins --- so the map's content lives in \emph{coordinates}
(where a patient sits) and in \emph{directions of extremity} (which corner they lean toward), not in
\emph{labels}. Any hard partition imposed on it is a convention, useful only insofar as a downstream task
rewards it. This is why the program never crowns a single number of clusters; it is also, as link three
will show, exactly the property that makes the map's signal a continuous gradient rather than a switch.
\end{definitionbox}

A continuum still has shape. Archetypal analysis reads that shape not as clusters but as \emph{corners} ---
the extreme profiles that the cloud is a blend of, with most patients ($61\%$) sitting in the interior as
mixtures rather than at any vertex. At the stable $A=4$ resolution (cross-seed Tucker congruence $0.997$) the
four corners are quasi-independent axes of extremity: $A_0$ a \emph{biological} corner (metabolic $+1.99$,
inflammatory $+2.52$, substance $+2.55$), $A_1$ a \emph{low-burden} pole, $A_2$ a \emph{severe but
non-biological} corner, $A_3$ a \emph{psychiatric-symptom} corner. The headline of the chapter is that
these three things --- biology, symptoms, severity --- vary nearly independently: a patient can be
biologically extreme without being symptomatically or globally severe, and vice versa. And of the four
corners, the biological one is the most distinct: it is the extreme that the symptom and severity axes
\emph{cannot} see, the corner you would miss entirely if you looked only at how ill a patient appears. Hold
that thought; the next two links are about that specific corner.

\subsection{Link two --- M3: that biological corner is the durable one}
The second link establishes \emph{time}. A map can be real and still be a snapshot of a passing state ---
in which case forecasting from it would be forecasting from noise. The temporal chapter (\S\ref{sec:temporal})
scored each follow-up visit onto the \emph{frozen} baseline map --- same loadings, same geometry, never
re-discovered --- and asked, axis by axis, how much of a patient's coordinate is a stable trait versus a
moving state. The instrument first had to survive the trip: measurement invariance holds, all five
continuous backbone axes reproducing across visits at $\varphi\geq0.95$ (metabolic $0.988$, inflammatory
$0.974$). Only once the ruler is shown not to warp can a coordinate change be read as the patient changing.

\begin{keyresult}[Durable biology, moving symptoms]
The test--retest reliabilities split the map cleanly along the very seam M2 found. The \emph{biological}
axes are traits: metabolic $\text{ICC}=0.91$, substance $0.88$, cognition $0.70$, inflammatory $0.62$ ---
a patient's biological position is close to a fixed property. The \emph{symptom} axes carry the state: sleep
($0.47$) and suicidality ($0.42$) are mixed, developmental ($0.39$) clearly state --- they slide visit to
visit. Severity is a trait by
rank ($0.62$) on a population that improves overall. Read against link one, the corner that M2 singled out
as most distinct (the biological one) is exactly the corner M3 finds most durable. Clinical translation,
carried forward verbatim: \emph{stratify on the durable biology, monitor the moving symptoms}.
\end{keyresult}

The two routes to this conclusion --- a variance decomposition (how much of each coordinate is trait) and a
geometry check (does a patient keep their archetype identity over time) --- agree, which is the kind of
internal corroboration that makes a finding hard to dismiss. Archetype identity persists across visits
(weight-cosine median $0.90$; $58\%$ of trajectories stable), and the two routes correlate ($\rho=-0.27$):
patients whose coordinates are most trait-like are the ones whose archetype membership is most stable. One
honest seam in this link, flagged here and revisited in the limitations: the \emph{geometry} route
(reliable-change on archetype weights) is precision-confounded --- patients measured more precisely look
more stable for a reason that has nothing to do with biology --- so the clean trait/state signal is the
error-corrected variance route, the ICCs above, not the geometry overlay. The biology-is-durable claim
rests on the clean route.

\subsection{Link three --- M4: the durable biological corner predicts functioning}
The third link converts durability into a forecast, and it is the link the program was built to test:
\emph{persists} is a property of the past, \emph{predicts} is a claim about the future, and they are not the
same thing. The prognosis chapter (\S\ref{sec:prognosis}) set a deliberately hard bar --- the map is
credited only if a baseline coordinate forecasts a two-year outcome \emph{incrementally beyond} what a
clinician already holds: the DSM-5 diagnosis, the current severity, and the present value of the outcome
itself. Against that bar, held out, the archetype representation buys a real increment on two-year
\emph{functioning} (egf, $N=2{,}114$): $\Delta\text{ELPD}=+59.4\pm11.1$ for the $A=4$ archetypes, and it is
robust --- it survives attrition weighting, permutation, and dropping cohorts (with the caveat that the
signal is course-dependent, carried by the episodic bipolar and depression patients more than by the
baseline-saturated schizophrenia ones). Severity, by contrast, is autoregression-saturated: the future of
\emph{how ill} is mostly told by the present of how ill, and the map adds little. The map predicts the
\emph{trajectory of function}, not the level of severity.

\begin{keyresult}[The biological corner has the worst functional prognosis]
The prognostic atlas reads the forecast off the corners of link one. Ordered by two-year functional
remission (GAF $\geq 71$): the \emph{low-burden} pole $A_1$ remits at $60\%$, the symptom corner $A_3$ at
$43\%$, the severe-non-biological corner $A_2$ at $32\%$, and the \emph{biological corner} $A_0$ at $27\%$
--- the worst. The corner M2 flagged as most distinct and M3 found most durable is, in M4, the corner with
the dimmest functional outlook. That is the chain closing on itself: a real biological extreme, that lasts,
that predicts where functioning goes.
\end{keyresult}

Two honesties keep this link calibrated, and both are stated here rather than left to the reader to infer.
First, where the signal lives. The predictive object is the \emph{fuller $A=4$ archetype representation}
--- biology read as the $A_0$ corner of the whole simplex --- not an isolated three-axis ``durable block''
pulled out on its own (added alone, the three trait axes give only an ambiguous $+2.2\pm3.1$). The corner
is biological, but it earns its forecast as a \emph{phenotype} --- a joint position in the space --- not as
a lone biomarker. Second, the magnitude. A held-out evidence gain of $\Delta\text{ELPD}=+59$ is a
group-level result; pushed down to an individual yes/no call it collapses to $\Delta\text{AUC}=+0.011$
(functional-remission $0.745\to0.756$). The chain is genuine and it is modest: it sharpens \emph{who, on
average, is heading where}, not \emph{will this person remit}.

\begin{caveatbox}[What the chain does and does not buy]
The three links establish a \emph{scientifically valid} object: a real continuum whose biological corner is
distinct (M2), durable (M3), and incrementally predictive of two-year functioning, co-informative with
DSM-5 and robust (M4). They do \emph{not} establish strong clinical utility: the individual increment is
small ($\Delta\text{AUC}=+0.011$), the predictive content is phenotype-level rather than a single isolated
axis, the signal is course-dependent, and --- as the treatment chapter (\S\ref{sec:treatment}) shows --- on
observational treatment-as-usual the map does not reliably tell you \emph{what to do} about any of it. The
chain says \emph{stratify on the durable biological phenotype --- the archetype simplex --- and forecast
functioning at the group level}; it stops
short of an individual risk calculator or a prescribing rule. The full account of both bars is the subject
of the conclusion and the limitations.
\end{caveatbox}
